\section{Introduction}\label{sec:intro}

Class imbalance is a persistent challenge in machine learning. In production
settings the cost of the resulting performance degradation is not abstract: in
medical screening at 10\% prevalence, a one-percentage-point drop in specificity
at fixed sensitivity multiplies the false-positive volume that clinicians must
reconcile, and in fraud detection it directly inflates the analyst caseload that
determines operational feasibility. When one class dominates the training
distribution, classifiers tend to optimize for majority-class accuracy at the
expense of minority-class recall \cite{he2009learning}. SMOTE
\cite{chawla2002smote}, which generates synthetic minority instances by
interpolating between existing samples and their nearest neighbours, remains the
dominant solution after two decades and has spawned hundreds of variants
\cite{fernandez2018smote,kovacs2019smote}. Yet a stubborn regularity persists:
SMOTE variants almost invariably degrade majority-class accuracy while improving
minority-class recall. This accuracy--balanced-accuracy tradeoff is so common
that practitioners treat it as inherent to oversampling.

The field has responded in two ways. One line asks \emph{which} resampling
method a given dataset needs: data-complexity studies show that data
characteristics, not the imbalance ratio alone, moderate oversampling
effectiveness \cite{komorniczak2022data}, and meta-learning research recommends
a resampler from dataset meta-features
\cite{carvalho2025resampling,jiang2026beyond}. Another line has grown openly
critical: large benchmark studies find that oversampling often fails to
outperform doing nothing once evaluation is done carefully, and that reported
gains depend heavily on which metrics are chosen
\cite{tarawneh2022stop,elor2022smote}. Both lines are
\emph{correlational}: they relate measured data characteristics or benchmark
outcomes to observed performance, but do not explain \emph{why} a method helps
or harms, and the proposed selection rules barely beat the trivial ``always
SMOTE'' baseline. What is missing is a \emph{mechanistic} account of the
tradeoff that tells us which lever to pull.

We provide such a mechanistic account. We argue the tradeoff is \emph{not}
inherent to oversampling but
is a mechanistic consequence of \emph{where} SMOTE generates synthetic data.
Standard SMOTE uses every minority instance as an interpolation seed, including
those in regions of genuine class overlap near the decision boundary. Generating
there pushes the learned boundary into majority-class territory, recovering some
minority instances while misclassifying previously correct majority ones. We
establish this in two steps. On controlled Gaussian mixtures with a \emph{known}
Bayes boundary, standard SMOTE demonstrably places synthetic points \emph{across}
that boundary. On real data, where the boundary is unobserved, a controlled
interventional experiment manipulates \emph{where} synthetic points are generated:
augmenting \emph{only} the model-estimated boundary (irreducible) errors reproduces
the full SMOTE signature (balanced accuracy $p<0.001$, negative accuracy trend),
whereas augmenting non-boundary or random instances has no such effect. The
evidence indicates that much of the cost traces to \emph{where} synthetic data is
placed, rather than to rebalancing per se.

To localize this boundary/overlap region on real data we introduce \textbf{Error
Instance Triage}, which categorizes each misclassified training instance from
\emph{out-of-bag} random-forest confidence and local class representation; the
categories are interpreted through, and validated against, an
aleatoric--epistemic uncertainty decomposition
\cite{shaker2020aleatoric,hullermeier2021aleatoric}.
We state the scope plainly: the triage \emph{instrument} is computed from
random-forest ensembles (the standard vehicle for this decomposition
\cite{shaker2020aleatoric}), while the downstream conclusions are evaluated
across three base learners, gradient-boosted trees, random forests, and
logistic regression (\S\ref{sec:results:frontier}). The categories are:

\begin{enumerate}
\item \textbf{Cat\,1 (Noise):} likely mislabeled instances: high total but
  low aleatoric uncertainty; the label conflicts with the feature pattern.
\item \textbf{Cat\,2 (Data-limited):} errors in regions lacking same-class
  evidence: high epistemic uncertainty, low local class density. These are
  \emph{reducible}: more correct-class data would resolve them.
\item \textbf{Cat\,3 (Irreducible):} errors in estimated boundary/overlap regions
  where classes overlap: high aleatoric uncertainty. These are \emph{not}
  reducible by adding data.
\end{enumerate}

The decomposition yields a diagnostic, not a new oversampler. On real imbalanced
data the misclassified minority concentrates in these estimated boundary/overlap
regions (Cat\,3 is $\sim$85\% of errors, a fraction that is stable under strict
out-of-fold re-estimation), exactly where SMOTE's harm originates. Yet that
balanced-accuracy deficit turns out to be largely a \emph{decision-threshold
artefact}: moving the decision threshold alone (a non-generative change)
recovers roughly nine percentage points of balanced accuracy over the default
threshold (\S\ref{sec:results:frontier}), without any synthetic generation. The practical implication is that the lever is
the \emph{decision rule}, not synthetic generation, and indeed non-generative
strategies (threshold-moving, cost-sensitive learning, balanced ensembles)
capture the recoverable balanced accuracy \emph{without} the boundary-generation
accuracy cost.

We make this precise with a \emph{remedy decomposition} of the misclassified
minority that sorts each error into threshold-recoverable, data-reducible
(epistemic), or irreducible (aleatoric). We find the deficit is overwhelmingly
\textbf{threshold-recoverable} (mean $68\%$ across 79 datasets, $80\%$ at
$\mathrm{IR}>3$): the classifier's scores already place these instances above a
tuned operating point. The known
``move the threshold'' equivalence
\cite{provost2000imbalanced,elkan2001foundations,elor2022smote,goorbergh2022harm}
is the consequence: across a 16-strategy menu oversampling yields no ranking gain
and shifts the class prior in probability space rather than improving
calibration, and once \emph{every} strategy is given the same
out-of-fold threshold tuning the large default-threshold gap between the families
collapses to within $\pm1$\,pp across three base learners.
We turn the diagnosis into a prescription with a regime map and an
honestly-scoped selector, which beats the always-SMOTE / imbalance-ratio
paradigm, but only because the menu now includes the dominant non-generative
strategies, while a strong non-generative default is itself near-optimal.

\paragraph{Contributions.}
\begin{enumerate}
\item \textbf{A remedy decomposition of the minority deficit (headline).} We
  decompose the \emph{misclassified} minority by combining an error triage
  (out-of-bag confidence and local representation, validated against an
  aleatoric--epistemic uncertainty decomposition) with a threshold-recoverability
  test, mapping each error to the remedy that would fix it:
  \emph{threshold-recoverable}, \emph{data-reducible} (epistemic), or
  \emph{irreducible} (aleatoric). Unlike prior
  instance taxonomies, which score \emph{difficulty}
  \cite{napierala2016types,smith2014instance,swayamdipta2020dataset,seedat2022dataiq},
  ours conditions on the model being wrong and introduces the threshold-recoverable
  axis. Empirically the deficit is overwhelmingly threshold-recoverable
  ($68\%$ on average, $80\%$ at $\mathrm{IR}>3$).
\item \textbf{A mechanism behind the ``move the threshold'' equivalence.} The
  decomposition \emph{explains} a known regularity (that resampling does not
  improve ranking and a threshold shift reproduces its benefit
  \cite{provost2000imbalanced,elkan2001foundations,elor2022smote,goorbergh2022harm})
  by showing the deficit is mostly a thresholding artefact, so oversampling is
  an implicit, inefficient enactment of a threshold move. We give the geometric
  mechanism: on controlled mixtures with a \emph{known} boundary, SMOTE generates
  across it, and a short proposition (\S\ref{sec:method:theory}) shows why
  interpolating boundary seeds crosses into the majority-posterior region.
\item \textbf{Evidence at scale (confirmation, not claim).} Across 79 datasets and
  a 16-strategy menu we reproduce, on a broad roster, that oversampling yields no
  discrimination gain (AUC unchanged) and enacts a prior shift in probability
  space rather than a calibration improvement, and that once
  \emph{every} strategy receives the same out-of-fold threshold tuning, the
  $+5.76$\,pp default-threshold gap between the families collapses to within
  $\pm1$\,pp across three base learners (six evaluated; the documented exception
  is the single MLP, whose \emph{ranking} imbalance corrupts and resampling
  partially repairs, \S\ref{sec:results:frontier}). The apparent dominance of non-generative
  methods at the default threshold is an operating-point effect. We frame these as confirmation of
  \cite{elor2022smote,goorbergh2022harm,dalpozzolo2015calibrating} at scale.
\item \textbf{An honestly-scoped prescription.} The decomposition says which remedy
  applies; we add a regime map and, as a secondary analysis, a held-out selector
  that beats always-SMOTE by $3.5$ points, but an imbalance-ratio-only rule
  matches it and a strong non-generative default is near-optimal, so the
  \emph{menu}, not selector sophistication, is the lever (the practical answer to
  the selection problem left open by \cite{jiang2026beyond}).
\end{enumerate}
